\section*{Conclusion}
\addcontentsline{toc}{section}{Conclusion}

L'étude \acrshort{CFD} menée sur la fusée SP-02 avait pour objectif de valider théoriquement l'efficacité du système d'aérofreins
du premier étage. Les simulations réalisées sur six configurations d'ouverture montrent une augmentation monotone du coefficient
de traînée, de $0.5936$ aérofreins fermés à $0.7827$ en configuration totalement déployée, soit un gain de près de $32~\%$. Cette
évolution est correctement décrite par la loi quadratique établie à la section \ref{sec:cd}, avec un coefficient de détermination
de $0.999$.\\

L'analyse des champs de pression et de vitesse confirme par ailleurs la cohérence physique de ces résultats. Le déploiement des
aérofreins fait apparaître une zone de surpression sur leur face amont et une zone de dépression en aval, associée à un
décollement précoce de la couche limite et à un élargissement du sillage. Ce déséquilibre de pression constitue bien le mécanisme
dominant de génération de traînée, et son intensité croît avec le taux d'ouverture. En dehors des zones directement influencées
par les aérofreins, l'écoulement conserve un comportement classique, ce qui traduit la validité générale du modèle.\\

Le concept est donc validé sur le plan aérodynamique : le bloc aérofreins est capable de produire l'augmentation de traînée
nécessaire à la séparation des étages, et la loi obtenue fournit une donnée d'entrée directement exploitable pour la simulation
de vol et pour le dimensionnement des efforts repris par le mécanisme de déploiement.\\

Ces conclusions doivent toutefois être lues à la lumière des hypothèses retenues. La résolution stationnaire des équations
moyennées de Reynolds ne rend pas compte des fluctuations temporelles de l'écoulement, et le modèle $k-\varepsilon$ réalisable,
s'il constitue un bon compromis pour une étude comparative, reste approximatif dans les zones de fort décollement. Les aérofreins
ont de plus été simulés en positions figées et à incidence nulle, alors que la fusée présente en vol une incidence non nulle et
une phase transitoire de déploiement. Enfin, aucune étude de convergence en maillage ni aucune confrontation à des données
expérimentales n'a pu être conduite dans le cadre de ce projet.\\

Les valeurs absolues de $C_D$ sont donc à considérer avec la réserve qui s'impose. En revanche, la tendance dégagée et les ordres
de grandeur obtenus sont suffisamment robustes pour justifier la pertinence du concept et orienter la conception du bloc
aérofreins.
